\chapter{Problemas técnicos y soluciones}
\label{chap:problemas}

Este capítulo reúne los problemas más costosos del desarrollo. Para cada uno se describen el
síntoma, la causa y la solución, tal como quedaron documentados en el código y en el repositorio.
Muchos no se deben a un error de programación, sino a detalles de los componentes o del montaje que
no son evidentes en una primera lectura.

\section{UART sobre PIO y latencia del enlace Modbus}
\label{sec:prob-uart}

\paragraph{Síntoma.} Tras migrar el firmware del ESP32-S3 a la Pico 2, las órdenes tardaban entre 3 y
\SI{4}{\second} en tener efecto.

\paragraph{Causa.} En el ESP32-S3 la UART es un periférico hardware: los bytes de la respuesta llegan
a una FIFO por interrupción y, cuando ModbusMaster los pide, ya están disponibles, así que cada
transacción dura 1 a \SI{3}{\milli\second}. En la Pico, el trazado lleva la transmisión a GP14 y la
recepción a GP13, y se usó \codigo{SerialPIO}, una UART implementada con una máquina de estado PIO
que funciona en cualquier pin \verificar{confirmar si el RP2350 permite la función UART en GP14
mediante su función alternativa}. Con \codigo{SerialPIO}, el bucle de espera de ModbusMaster agotaba
su tiempo límite, de unos \SI{20}{\milli\second}, en muchas transacciones: diez escrituras
consecutivas bloqueaban el núcleo \SI{200}{\milli\second}, y con las lecturas de supervisión se
llegaba a varios segundos.

Un segundo detalle afectaba al cambio de sentido del transceptor: llamar a \codigo{flush()} antes de
desactivar el emisor, como se haría con una UART hardware, rompía la recepción inmediata de la
respuesta. La solución fue esperar un tiempo fijo de \SI{200}{\micro\second} tras el último byte.

\paragraph{Primera solución: doble núcleo.} La versión \fichero{firmwarerasp} dedicó el núcleo 1 al
bus Modbus, con la consigna publicada en variables \codigo{volatile} y un protocolo de relevo para
que el núcleo 0 pudiera tomar el bus. Funcionaba, pero con variables compartidas sin atomicidad,
esperas de relevo y lecturas cuyo instante no estaba ligado al de la muestra de telemetría.

\paragraph{Solución final: dueño único con lazo de periodo fijo.} Ajustados los tiempos de espera y
el cambio de sentido, cada transacción volvió a durar pocos milisegundos y se pudo prescindir del
segundo núcleo: un único lazo de \SI{50}{\milli\second} con exactamente tres transacciones por
vuelta. El resultado es un tiempo de trabajo de \SI{11.0}{\milli\second} con \SI{0.2}{\milli\second}
de dispersión (sección~\ref{sec:res-temporizacion}), sin condiciones de carrera y con lecturas,
lógica y escritura en la misma vuelta.

\section{Configuración del ZSC31014: la ventana del modo comando}
\label{sec:prob-zsc}

\paragraph{Contexto.} Para cambiar la ganancia o el cero del acondicionador hay que escribir en su
EEPROM, y para ello entrar en modo comando. Según la hoja de datos \cite{zsc31014}, tras un reinicio
por encendido (POR) el chip espera la orden Start\_CM durante \SI{6}{\milli\second} si la EEPROM no
está bloqueada, o \SI{1.5}{\milli\second} si lo está. Pasada la ventana, ignora Start\_CM sin dar
ninguna señal de error.

\subsection{Primer obstáculo: provocar el reinicio}

Tras reprogramar la Pico, la configuración nunca se aplicaba, porque la placa Click seguía
alimentada y sin POR no se abre la ventana. La solución fue controlar su alimentación con el pin EN
(GPIO1) y hacer un ciclo completo de apagado y encendido antes de cada intento.

\subsection{Segundo obstáculo: alimentación parásita por el bus}

\paragraph{Síntoma.} Con EN a nivel bajo, el chip seguía respondiendo por I\textsuperscript{2}C, y la
tensión de la Click se quedaba en \SIrange{2.5}{2.7}{\volt} en lugar de caer a cero.

\paragraph{Causa.} Las resistencias de \emph{pull-up} de SDA, SCL e INT cuelgan del
\SI{3.3}{\volt} permanente del zócalo mikroBUS, no del carril gobernado por EN. Con el chip sin
alimentar, la corriente entraba por esos pines, atravesaba los diodos de protección internos y
alimentaba el carril. La hoja de datos sitúa el umbral de POR entre \num{1.8} y \SI{2.5}{\volt}:
con \SIrange{2.5}{2.7}{\volt} parásitos, el chip nunca veía un reinicio.

\paragraph{Diagnóstico y solución.} Se añadió la orden LC\_HOLD, que mantiene la célula sin
alimentación \SI{8}{\second} sin bloquear el lazo, con dos modos para las líneas del bus. Las
medidas con polímetro mostraron que sólo forzándolas a nivel bajo la tensión cae a cero, y así se
dejó por defecto. Hubo además un error en el propio diagnóstico: la primera versión de la prueba
dejaba las líneas en alto mientras medía, de modo que siempre ``detectaba'' realimentación y acusaba
al pin EN sin motivo.

\subsection{Tercer obstáculo: acertar la ventana}

La orden tiene que llegar después de que el chip arranque y antes de que la ventana se cierre, y ese
instante depende de los condensadores de la placa Click, así que los retardos copiados de otros
controladores no servían. En lugar de adivinarlo, el firmware lo mide: sondea la dirección hasta que
el chip responde, y a partir de ahí prueba retardos de \num{0} a \SI{1050}{\micro\second} en pasos
de \SI{150}{\micro\second} a lo largo de hasta ocho ciclos de alimentación, confirmando la entrada
con la lectura de B\_Config. El intento que funciona se devuelve a la interfaz para poder fijarlo
más adelante.

\subsection{Cuarto obstáculo: integridad de la EEPROM}

La firma MISR que protege la EEPROM se recalcula al salir del modo comando con Start\_NOM y se
comprueba en el siguiente POR: cortar la alimentación antes deja la firma a medias y el chip arranca
desde entonces en estado de diagnóstico. Por eso, tras cualquier escritura se envía Start\_NOM y se
espera diez veces el tiempo típico antes de cortar. Como ZMDI\_Config1 sólo se carga tras un POR, la
verificación se hace en una segunda pasada desde un reinicio limpio.

\section{Lecturas de la célula siempre ``antiguas''}

En una versión del firmware no se registraba ni una sola muestra válida. La lectura se hacía en dos
fases separadas \SI{2}{\milli\second}, quedándose con la segunda; como en modo de conversión
continua leer no dispara una conversión y el chip no había producido un dato nuevo en ese tiempo, la
segunda lectura venía siempre marcada como repetida. Una versión anterior funcionaba por casualidad,
porque sus dos lecturas caían en vueltas distintas del lazo. La solución fue una única lectura por
vuelta, aceptada sólo si el dato es nuevo.

\section{Depuración del ciclo de ensayo}
\label{sec:prob-ciclo}

Al reorganizar el ciclo para que empezara en el extremo inferior (sección~\ref{sec:fw-ciclo})
aparecieron tres problemas encadenados:

\begin{description}[style=nextline,leftmargin=1.5em]
  \item[El ciclo terminaba sin hacer nada] Con el eje ya en el final de carrera, la condición de
        llegada se cumplía de inmediato y el ciclo recorría todas sus fases en un par de periodos.
        No era un fallo de lógica, sino de arranque: faltaba la fase de colocación en B.
  \item[No quedaba rastro en el registro] Un ciclo que empieza y acaba en \SI{200}{\milli\second} no
        aparecía en el registro, porque sólo se escribía una línea por lote de telemetría y ninguna
        al cambiar de fase. Ahora cada cambio de fase escribe su propia línea, y la línea de
        telemetría incluye posición, finales de carrera y fase.
  \item[La interfaz no mostraba la fase] El firmware había pasado a enviar una trama de ciclo de 18
        bytes y el servidor, que esperaba 14, la descartaba entera. Se añadió la compatibilidad con
        la trama antigua para que un desajuste de versiones degrade la información en lugar de
        hacerla desaparecer.
\end{description}

Además, el ciclo pasó a reconocer los extremos por los finales de carrera en vez de por la posición
integrada. Con la posición, cualquier error de escala del accionamiento o cualquier cero mal
colocado desplazaba el recorrido; con los finales, el ensayo recorre siempre lo mismo.

\section{El modelo que no cerraba: la salida del cable desplazada}
\label{sec:prob-dx}

\paragraph{Síntoma.} Con el primer modelo, que situaba la salida del cable justo encima de la
articulación, la curva teórica no encajaba con los ensayos: el recorrido angular que predecía no era
el que hacía la barra, y el ajuste sólo cuadraba forzando valores poco creíbles de la reducción.

\paragraph{Causa.} La salida del cable no está sobre la vertical de la articulación, sino desplazada
horizontalmente. Ese desplazamiento no es un detalle de dibujo: con la salida en la vertical, el
coseno del ángulo de la barra se cancela con el seno del ángulo del cable y la tensión resulta
proporcional a la longitud de cable, que es el resultado elegante de la
ecuación~\eqref{eq:tension0}. Con la salida desplazada esa cancelación desaparece
(sección~\ref{sec:mod-tension}).

\paragraph{Solución.} Se reescribió el modelo con el desplazamiento $dx$ como parámetro y con el
ángulo medido entre la línea que une articulación y salida y la propia barra, que además elimina la
ambigüedad que tenía medir desde la horizontal. El modelo pasó a estar implementado una sola vez y a
alimentarse de una configuración guardada en el servidor.

\section{Lo que la fuerza no puede decidir}
\label{sec:prob-signo}

Al ajustar los parámetros del mecanismo contra la señal de la célula se obtuvieron coeficientes de
determinación altos ($R^2 \approx 0{,}93$ en los ensayos limpios), lo que parecía confirmar la
geometría. No lo confirmaba. Como la célula no está calibrada, el ajuste incluye la sensibilidad
como incógnita, y $R^2$ es el cuadrado de la correlación: no distingue una sensibilidad positiva de
una negativa. Dos geometrías que predicen tensiones con tendencias opuestas dan exactamente el mismo
$R^2$ (sección~\ref{sec:mod-ambiguedad}).

La consecuencia práctica llegó al añadir el par del accionamiento al contraste: con los mismos
parámetros, el par medido y el del modelo resultaron \emph{anticorrelados}
(capítulo~\ref{chap:resultados}). El análisis se reorganizó entonces alrededor del par, que no
depende de ninguna calibración pendiente, y la señal de la célula quedó como comprobación
secundaria.

\section{Tramo inicial sin carga y calibraciones incoherentes}

Los ensayos de la segunda campaña muestran que la primera media revolución de motor transcurre sin
tensión en el cable y con un par casi nulo: el mecanismo recoge holgura antes de empezar a levantar
la masa. No es un fallo, pero invalida ese tramo para el análisis y obliga a definir explícitamente
el recorrido con carga.

Por otra parte, la constante que relaciona cuentas y newtons ha resultado distinta en cada intento:
\num{14.41} cuentas/N en la configuración de la interfaz, unas \num{20} al ajustar contra el
modelo y \SIrange{27}{28}{} en los ensayos limpios; y en la campaña del 15 de septiembre el ajuste
salía incluso con el signo contrario. Todo ello confirma que la célula necesita una calibración con
masas patrón antes de dar valores absolutos de fuerza.

\section{Segunda célula de carga en el mismo bus}
\label{sec:prob-nau}

Se evaluó añadir una segunda célula con una Load Cell 4 Click (NAU7802) en el mismo bus. Surgieron
dos problemas. El primero, de convivencia: el ZSC31014 no toleraba transacciones ajenas entre la
petición y la recogida del dato, y de ahí nació el árbitro de \fichero{i2c\_bus.h}. El segundo, de
ruido: con ganancia 128 y en reposo se midieron unas 18.000 cuentas pico a pico con el regulador
interno a \SI{4.5}{\volt}, unas 3.500 bajándolo a \SI{3.0}{\volt} con calibración interna de cero, y
unas 350 al activar el \emph{chopper}, equivalentes a \SI{0.49}{\micro\volt} en la entrada. La
segunda célula se retiró del firmware final.

\section{Discrepancias con la hoja de datos}
\label{sec:prob-estado}

Al contrastar el firmware con la hoja de datos del ZSC31014 se encontraron dos discrepancias que no
afectan a los ensayos, pero conviene corregir:

\begin{enumerate}
  \item \textbf{Bits de estado.} Según la hoja de datos, 00 = dato válido, 01 = modo comando,
        10 = dato repetido y 11 = diagnóstico. El firmware define \codigo{LC\_STATUS\_STALE = 1} y
        \codigo{LC\_STATUS\_COMMAND = 2}, intercambiados. En la lectura normal sólo se acepta 00, así
        que no hay efecto; en la confirmación del modo comando la respuesta real (\codigo{0x5A}) se
        acepta por la segunda condición del código.
  \item \textbf{Codificación de la ganancia.} El campo PreAmp\_Gain no es lineal: 000 = ×1,5,
        100 = ×3, 001 = ×6, 101 = ×12, 010 = ×24, 110 = ×48, 011 = ×96 y 111 = ×192, es decir, los
        tres bits en orden inverso. El selector de la interfaz supone un código lineal, de modo que
        para los códigos 1 a 6 la ganancia grabada no es la mostrada. Los códigos 0 y 7 coinciden.
        \verificar{comprobar qué código de ganancia estaba grabado durante los ensayos}
\end{enumerate}

\section{Otros problemas}

\begin{description}[style=nextline,leftmargin=1.5em]
  \item[Órdenes perdidas por los atajos de teclado] Un byte de CRC igual a una letra de atajo hacía
        que se perdiera la orden y se reiniciara el servo. Se resolvió aceptando los atajos sólo con
        el decodificador en reposo.
  \item[Frecuencia de muestreo imposible] La interfaz llegó a mostrar \SI{3813}{\hertz}: tras un
        reinicio del microcontrolador su reloj volvía a cero y la ventana de medida no se podaba
        nunca. Se detecta el salto de reloj y se reinicia la ventana.
  \item[Desplazamientos fantasma] Tras la programación de la célula, que bloquea el lazo cerca de un
        segundo, la integración multiplicaba ese hueco por la última velocidad. Se anula la base de
        tiempo tras cada operación bloqueante.
  \item[Órdenes que ``no hacían nada''] Una acción mal escrita en la interfaz se ignoraba sin avisar.
        Se añadió una confirmación para toda orden, con el motivo del fallo cuando lo hay.
  \item[Revisión A de la placa] Conflictos de red en el bloque RS-485, entradas analógicas en pines
        sin ADC y amplificadores sin alimentar (tabla~\ref{tab:erratas}).
\end{description}

\section{Lecciones aprendidas}

\begin{itemize}
  \item \textbf{Leer la hoja de datos entera}: la ventana del modo comando, la firma de la EEPROM, el
        umbral de POR y la codificación no lineal de la ganancia están en secciones distintas y todas
        han resultado críticas.
  \item \textbf{Medir antes de suponer}: la alimentación parásita sólo se descubrió con un polímetro,
        y el retardo de la ventana, con un barrido.
  \item \textbf{Un ajuste bueno no es una validación}: un $R^2$ alto puede esconder una ambigüedad de
        signo. Hace falta una segunda magnitud, medida de forma independiente, para cerrar el
        contraste.
  \item \textbf{Cada operación debe responder}: las confirmaciones de órdenes y el registro por
        cambio de fase convirtieron fallos silenciosos en diagnósticos claros.
  \item \textbf{Conservar los datos brutos}: ha permitido recalibrar y reanalizar todos los ensayos
        sin repetirlos, incluso después de cambiar el modelo por completo.
\end{itemize}
